\scanpage{395}\mspage{\unsure{119|149}}
alors la flèche de chgt. de base
\[
p^{\ast}f_{!}\longrightarrow f'_{!}\unsure{g|q}^{\ast}
\]
entre foncteurs $\mathrm{HOT}_{W}(X)\longrightarrow\mathrm{HOT}_{W}(S')$ est iso.

C'est pratiquement trivial, compte tenu, dans
le cas $f$ $W$-lisse, que si $Z$ sur $X$ est \unsure{$W$-lisse|lisse} sur $X$,
il l'est sur $S$ ; et dans le cas $p$ $W$-propre, en
utilisant le dernier énoncé (``oublié'') du (6),
p.~199.

7.8. Je termine en mettant en vedette
ce qui, à présent, me paraît la
\deleted{\unsure{résultat}} \inserted{substance} principale de tous ces ``rappels'', et
comme la quintessence de la théorie des localiseurs
fondamentaux développée jusqu'à présent :
\[
\left|\left\Vert\left\{\begin{aligned}
\mathrm{Loc}(3\ \mathrm{bis},)\deleted{\unsure{6\mathrm{a})}}&\Longrightarrow\mathrm{Der}(4^{\unsure{\alpha|\prime}},5^{\alpha})\\
\mathrm{Loc}(3\ \mathrm{bis},6\mathrm{a})&\Longrightarrow\mathrm{Der}(1,2,3,4^{\alpha},5^{\alpha},6^{\alpha})\\
\mathrm{Loc}(6'\mathrm{a},8,9)&\Longrightarrow\text{tous les axiomes des}\\
&\hphantom{\Longrightarrow{}}\text{dérivateurs}\quad(+\mathrm{Loc}(1\ \mathrm{ter})\deleted{3\ \mathrm{bis})})
\end{aligned}\right.\right.\right.
\]

Je considère les axiomes $\mathrm{Loc}(3\ \mathrm{bis})$,
$\mathrm{Loc}(8)$, $\mathrm{Loc}(9)$, comme des axiomes
``anodins'' --- tous les localiseurs \unsure{qu'on|qui}
\unsure{définit|définit.} par voie cohomologique ou homotopique
ordinaire, y satisfont. Il reste l'axiome
$\mathrm{Loc}(6\mathrm{a})$ \deleted{\unsure{des caté}} \inserted{ou $\mathrm{Loc}(6'\mathrm{a})$} \unsure{$W$-cofibrations|$W$-cofibration} (\unsure{formes|sens} faibles),
% archive sheet number: 394
